% Figure: power reflectance of s- and p-polarised light versus angle of
% incidence at an air-to-glass boundary (n1=1, n2=1.5), showing the
% Brewster zero of the p curve near 56 degrees and the rise of both to
% unity at grazing incidence. Fresnel intensity reflectances plotted with
% Snell's law folded in, cos(theta2) written through sin: with
% s2 = sin(th1)/1.5, cos(th2) = sqrt(1 - s2^2). pgfplots evaluates the
% degree-argument trig, so angles are passed in degrees and the sqrt
% argument stays non-negative across the whole 0..90 domain.
\begin{figure}[htbp]
\centering
\begin{tikzpicture}
\begin{axis}[
  width=12cm, height=6.8cm,
  xlabel={angle of incidence $\theta_1$ (degrees)},
  ylabel={power reflectance $R$},
  xmin=0, xmax=90, ymin=0, ymax=1,
  axis lines=left,
  legend pos=north west,
  legend cell align=left,
  domain=0:89.5, samples=300,
  xtick={0,15,30,45,60,75,90},
  ytick={0,0.2,0.4,0.6,0.8,1.0},
]
  % s-polarisation: r_s = (cos th1 - 1.5 cos th2)/(cos th1 + 1.5 cos th2),
  % with cos th2 = sqrt(1 - (sin th1 / 1.5)^2). R_s = r_s^2.
  \addplot[engblue, very thick] {
    ( ( cos(x) - 1.5*sqrt( 1 - (sin(x)/1.5)^2 ) )
    / ( cos(x) + 1.5*sqrt( 1 - (sin(x)/1.5)^2 ) ) )^2
  };
  \addlegendentry{$R_s$ (perpendicular)}
  % p-polarisation: r_p = (1.5 cos th1 - cos th2)/(1.5 cos th1 + cos th2).
  \addplot[engred, very thick] {
    ( ( 1.5*cos(x) - sqrt( 1 - (sin(x)/1.5)^2 ) )
    / ( 1.5*cos(x) + sqrt( 1 - (sin(x)/1.5)^2 ) ) )^2
  };
  \addlegendentry{$R_p$ (parallel)}
  % Brewster marker at arctan(1.5) = 56.31 deg, where R_p = 0.
  \addplot[enggray, dashed, thick] coordinates {(56.31,0) (56.31,1)};
  \node[enggray, font=\scriptsize, anchor=south, rotate=90]
    at (axis cs:56.31,0.30) {$\theta_B = 56.3^{\circ}$};
  \node[engred, font=\scriptsize, anchor=west]
    at (axis cs:44,0.05) {$R_p \to 0$};
\end{axis}
\end{tikzpicture}
\caption[Fresnel reflectance and Brewster's angle at an air-glass
boundary]{Power reflectance of the two polarisations against angle of
incidence for light striking glass ($n_1 = 1$, $n_2 = 1.5$) from air.
Both curves begin at the $4\,\%$ normal-incidence value and climb to
total reflection at grazing incidence, but the parallel ($p$) reflectance
dips to exactly zero at Brewster's angle $\theta_B = \arctan(n_2/n_1) =
56.3^{\circ}$, where the reflected ray is purely $s$-polarised. A window
tilted to this angle passes $p$-polarised light without a reflection
loss, the working geometry of a laser-cavity window and a polarising
beam pick-off. The perpendicular ($s$) reflectance rises monotonically
and never vanishes, which is why an unpolarised beam reflected near
Brewster's angle emerges strongly polarised.}
\label{fig:vol04ch10:brewster-reflectance}
\end{figure}
